% GENERATED FILE. Edit canonical lessons, then run npm run generate:books.
% canonical-source-hash: fnv1a64:00f1cc6516d45478
% canonical-lessons: TE-C119-tene, TE-C119-idli, TE-C119-dosa, TE-C119-vada, TE-C119-kaphi, TE-R119-first-pass-the-house-and-things-you-carry, TE-R119-second-pass-clothes-and-food

\chapter{Honey, Idli, Dosa, Vada, Coffee}
\label{ch:te-honey-idli-dosa-vada-coffee}
\hlchaptermodality{119}

\begin{chapteropening}
\textbf{By the end of this chapter:} \emph{I can say honey, an idli, a dosa, a vada and coffee.}

Complete the last lesson of chapter 119: i can say honey, an idli, a dosa, a vada and coffee.
\end{chapteropening}

\section[\texorpdfstring{\te{తేనె} (tēne)}{తేనె (tēne)}]{\te{తేనె} (tēne) — honey}

\label{lesson:TE-C119-tene}



\begin{quote}
Before the new one: say the Telugu for ginger, then the Telugu for garlic.
\end{quote}

\subsection*{You'll want to know: \te{తేనె}}
\textbf{\te{తేనె}} — \emph{tēne} — \textquotedblleft{}honey\textquotedblright{}.

\begin{grammarlens}[title={What you've built}]
One more word for this chapter.
\end{grammarlens}

\subsection*{Your turn}
\begin{itemize}
\raggedright
  \item \emph{Say it:} \emph{tēne}
  \item \emph{Say it:} \emph{tēne}, once more
  \item \emph{Say it:} say \emph{tēne}
  \item \emph{Recall:} say the Telugu for ginger, then the Telugu for garlic, then say \emph{tēne} again
\end{itemize}

\begin{tcolorbox}[breakable,colback=teal!4,colframe=teal!35!black,title={Before you move on}]
What does \te{తేనె} mean? (\textquotedblleft{}Honey\textquotedblright{}.) Say it once more.
\end{tcolorbox}
\section[\texorpdfstring{\te{ఇడ్లీ} (iḍlī)}{ఇడ్లీ (iḍlī)}]{\te{ఇడ్లీ} (iḍlī) — an idli}

\label{lesson:TE-C119-idli}



\begin{quote}
Before the new one: say the Telugu for garlic, then the Telugu for honey.
\end{quote}

\subsection*{You'll want to know: \te{ఇడ్లీ}}
\textbf{\te{ఇడ్లీ}} — \emph{iḍlī} — \textquotedblleft{}an idli\textquotedblright{}.

\begin{grammarlens}[title={What you've built}]
One more word for this chapter.
\end{grammarlens}

\subsection*{Your turn}
\begin{itemize}
\raggedright
  \item \emph{Say it:} \emph{iḍlī}
  \item \emph{Say it:} \emph{iḍlī}, once more
  \item \emph{Say it:} say \emph{iḍlī}
  \item \emph{Recall:} say the Telugu for garlic, then the Telugu for honey, then say \emph{iḍlī} again
\end{itemize}

\begin{tcolorbox}[breakable,colback=teal!4,colframe=teal!35!black,title={Before you move on}]
What does \te{ఇడ్లీ} mean? (\textquotedblleft{}An idli\textquotedblright{}.) Say it once more.
\end{tcolorbox}
\section[\texorpdfstring{\te{దోశ} (dōsa)}{దోశ (dōsa)}]{\te{దోశ} (dōsa) — a dosa}

\label{lesson:TE-C119-dosa}



\begin{quote}
Before the new one: say the Telugu for honey, then the Telugu for an idli.
\end{quote}

\subsection*{You'll want to know: \te{దోశ}}
\textbf{\te{దోశ}} — \emph{dōsa} — \textquotedblleft{}a dosa\textquotedblright{}.

\begin{grammarlens}[title={What you've built}]
One more word for this chapter.
\end{grammarlens}

\subsection*{Your turn}
\begin{itemize}
\raggedright
  \item \emph{Say it:} \emph{dōsa}
  \item \emph{Say it:} \emph{dōsa}, once more
  \item \emph{Say it:} say \emph{dōsa}
  \item \emph{Recall:} say the Telugu for honey, then the Telugu for an idli, then say \emph{dōsa} again
\end{itemize}

\begin{tcolorbox}[breakable,colback=teal!4,colframe=teal!35!black,title={Before you move on}]
What does \te{దోశ} mean? (\textquotedblleft{}A dosa\textquotedblright{}.) Say it once more.
\end{tcolorbox}
\section[\texorpdfstring{\te{వడ} (vaḍa)}{వడ (vaḍa)}]{\te{వడ} (vaḍa) — a vada}

\label{lesson:TE-C119-vada}



\begin{quote}
Before the new one: say the Telugu for an idli, then the Telugu for a dosa.
\end{quote}

\subsection*{You'll want to know: \te{వడ}}
\textbf{\te{వడ}} — \emph{vaḍa} — \textquotedblleft{}a vada\textquotedblright{}.

\begin{grammarlens}[title={What you've built}]
One more word for this chapter.
\end{grammarlens}

\subsection*{Your turn}
\begin{itemize}
\raggedright
  \item \emph{Say it:} \emph{vaḍa}
  \item \emph{Say it:} \emph{vaḍa}, once more
  \item \emph{Say it:} say \emph{vaḍa}
  \item \emph{Recall:} say the Telugu for an idli, then the Telugu for a dosa, then say \emph{vaḍa} again
\end{itemize}

\begin{tcolorbox}[breakable,colback=teal!4,colframe=teal!35!black,title={Before you move on}]
What does \te{వడ} mean? (\textquotedblleft{}A vada\textquotedblright{}.) Say it once more.
\end{tcolorbox}
\section[\texorpdfstring{\te{కాఫీ} (kāphī)}{కాఫీ (kāphī)}]{\te{కాఫీ} (kāphī) — coffee}

\label{lesson:TE-C119-kaphi}



\begin{quote}
Before the new one: say the Telugu for a dosa, then the Telugu for a vada.
\end{quote}

\subsection*{You'll want to know: \te{కాఫీ}}
\textbf{\te{కాఫీ}} — \emph{kāphī} — \textquotedblleft{}coffee\textquotedblright{}.

\begin{grammarlens}[title={What you've built}]
One more word for this chapter.
\end{grammarlens}

\subsection*{Your turn}
\begin{itemize}
\raggedright
  \item \emph{Say it:} \emph{kāphī}
  \item \emph{Say it:} \emph{kāphī}, once more
  \item \emph{Say it:} say \emph{kāphī}
  \item \emph{Recall:} say the Telugu for a dosa, then the Telugu for a vada, then say \emph{kāphī} again
\end{itemize}

\begin{tcolorbox}[breakable,colback=teal!4,colframe=teal!35!black,title={Before you move on}]
What does \te{కాఫీ} mean? (\textquotedblleft{}Coffee\textquotedblright{}.) Say it once more.
\end{tcolorbox}
\section[(dialogue)]{The house and things you carry — a first pass}

\label{lesson:TE-R119-first-pass-the-house-and-things-you-carry}



\begin{quote}
Say the words for a vada and coffee.
\end{quote}

\subsection*{Your turn}
Say these aloud:

\begin{itemize}
\raggedright
  \item a lock, a key, a rope, a nail, scissors
  \item a comb, soap, a towel, glasses, a clock
  \item a bag, a pen, paper, a drinking glass, a spoon
  \item a stove, a candle, a matchbox, a doll; a picture, a ball
  \item a kite, a shirt, a sari, a dhoti, sandals
\end{itemize}

\begin{tcolorbox}[breakable,colback=teal!4,colframe=teal!35!black,title={Before you move on}]
A lock? (\textbf{\te{తాళం}}, \emph{tāḷaṁ}.) A key? (\textbf{\te{తాళంచెవి}}, \emph{tāḷaṁcevi}.) A rope? (\textbf{\te{తాడు}}, \emph{tāḍu}.) A nail? (\textbf{\te{మేకు}}, \emph{mēku}.) Scissors? (\textbf{\te{కత్తెర}}, \emph{katterā}.) A comb? (\textbf{\te{దువ్వెన}}, \emph{duvvena}.) Soap? (\textbf{\te{సబ్బు}}, \emph{sabbu}.) A towel? (\textbf{\te{తువ్వాలు}}, \emph{tuvvālu}.) Glasses? (\textbf{\te{కళ్ళజోడు}}, \emph{kaḷḷajōḍu}.) A clock? (\textbf{\te{గడియారం}}, \emph{gaḍiyāraṁ}.) A bag? (\textbf{\te{సంచి}}, \emph{sañci}.) A pen? (\textbf{\te{కలం}}, \emph{kalaṁ}.) Paper? (\textbf{\te{కాగితం}}, \emph{kāgitaṁ}.) A drinking glass? (\textbf{\te{గ్లాసు}}, \emph{glāsu}.) A spoon? (\textbf{\te{చెంచా}}, \emph{cemcā}.) A stove? (\textbf{\te{పొయ్యి}}, \emph{poyyi}.) A candle? (\textbf{\te{కొవ్వొత్తి}}, \emph{kovvotti}.) A matchbox? (\textbf{\te{అగ్గిపెట్టె}}, \emph{aggipeṭṭe}.) A doll; a picture? (\textbf{\te{బొమ్మ}}, \emph{bomma}.) A ball? (\textbf{\te{బంతి}}, \emph{banti}.) A kite? (\textbf{\te{గాలిపటం}}, \emph{gālipaṭaṁ}.) A shirt? (\textbf{\te{చొక్కా}}, \emph{cokkā}.) A sari? (\textbf{\te{చీర}}, \emph{cīra}.) A dhoti? (\textbf{\te{పంచె}}, \emph{pañce}.) Sandals? (\textbf{\te{చెప్పులు}}, \emph{ceppulu}.)
\end{tcolorbox}
\section[(dialogue)]{Clothes and food — a second pass}

\label{lesson:TE-R119-second-pass-clothes-and-food}



\begin{quote}
Say the words for a cap and a ring.
\end{quote}

\subsection*{Your turn}
Say these aloud:

\begin{itemize}
\raggedright
  \item a cap, a ring, bangles, a curry, dal
  \item sambar, rasam, an egg, a fish, meat
  \item vegetables, a banana, a mango, coconut, an onion
  \item a tomato, a potato, an aubergine, ginger, garlic
  \item honey, an idli, a dosa, a vada, coffee
\end{itemize}

\begin{tcolorbox}[breakable,colback=teal!4,colframe=teal!35!black,title={Before you move on}]
A cap? (\textbf{\te{టోపీ}}, \emph{ṭōpī}.) A ring? (\textbf{\te{ఉంగరం}}, \emph{uṅgaraṁ}.) Bangles? (\textbf{\te{గాజులు}}, \emph{gājulu}.) A curry? (\textbf{\te{కూర}}, \emph{kūra}.) Dal? (\textbf{\te{పప్పు}}, \emph{pappu}.) Sambar? (\textbf{\te{సాంబారు}}, \emph{sāmbāru}.) Rasam? (\textbf{\te{రసం}}, \emph{rasaṁ}.) An egg? (\textbf{\te{గుడ్డు}}, \emph{guḍḍu}.) A fish? (\textbf{\te{చేప}}, \emph{cēpa}.) Meat? (\textbf{\te{మాంసం}}, \emph{māṁsaṁ}.) Vegetables? (\textbf{\te{కూరగాయలు}}, \emph{kūragāyalu}.) A banana? (\textbf{\te{అరటిపండు}}, \emph{araṭipaṇḍu}.) A mango? (\textbf{\te{మామిడిపండు}}, \emph{māmiḍipaṇḍu}.) Coconut? (\textbf{\te{కొబ్బరి}}, \emph{kobbari}.) An onion? (\textbf{\te{ఉల్లిపాయ}}, \emph{ullipāya}.) A tomato? (\textbf{\te{టమాటా}}, \emph{ṭamāṭā}.) A potato? (\textbf{\te{బంగాళాదుంప}}, \emph{baṅgāḷādumpa}.) An aubergine? (\textbf{\te{వంకాయ}}, \emph{vaṅkāya}.) Ginger? (\textbf{\te{అల్లం}}, \emph{allaṁ}.) Garlic? (\textbf{\te{వెల్లుల్లి}}, \emph{vellulli}.) Honey? (\textbf{\te{తేనె}}, \emph{tēne}.) An idli? (\textbf{\te{ఇడ్లీ}}, \emph{iḍlī}.) A dosa? (\textbf{\te{దోశ}}, \emph{dōsa}.) A vada? (\textbf{\te{వడ}}, \emph{vaḍa}.) Coffee? (\textbf{\te{కాఫీ}}, \emph{kāphī}.)
\end{tcolorbox}
